\documentclass[12pt]{article}
\usepackage[utf8]{inputenc}
\usepackage[T1]{fontenc}
\usepackage{hyperref}
\title{San Juan Teita Mixtec Language Resources for K-12 Curriculum Development}
\author{Compiled by Claude Sonnet 5.5 (claude-sonnet-5-5) for Andr\'as Kornai}
\date{}
\begin{document}
\maketitle
\section{Introduction}
San Juan Teita Mixtec (ISO 639-3 xtj, Glottocode sanj1282) is an Otomanguean language of Oaxaca, Mexico. Its endonym is reported as Da\~nudavi. Search results give roughly 370 to 400 speakers and describe it as vigorous (EGIDS 6a). Very little material exists, and the search found no school materials.
\section{Primary Educational Resources}
No primers, textbooks or school programmes were found in the search results.
\section{Linguistic and Reference Materials}
\begin{itemize}
\item ASJP database entry (word list data, speaker count): \url{https://asjp.clld.org/languages/SAN_JUAN_TEITA_MIXTEC}
\item Wikitongues language page: \url{https://wikitongues.org/languages/xtj/}
\end{itemize}
\section{Religious and Community Resources}
\begin{itemize}
\item Joshua Project people-group profile (San Juan Teita Mixteco): \url{https://joshuaproject.net/people_groups/14701/mx}
\end{itemize}
\section{Digital and Audio Resources}
\begin{itemize}
\item Global Recordings Network, Mixtec, San Juan Teita (audio and slideshow Bible lessons; some recordings may have been withdrawn): \url{https://globalrecordings.net/en/language/xtj}
\item Global Recordings Network, San Juan El Alto variety page: \url{https://glorecnet-website.k.globalrecordings.net/en/language/27261}
\end{itemize}
\section{Curriculum Development Context}
The language is small and little documented in what was found. A curriculum would need new materials developed with the community. The Global Recordings audio is evangelistic and not designed for school use.
\end{document}